\documentclass[11pt]{article}
\usepackage[margin=1.1in]{geometry}
\usepackage{graphicx,booktabs,amsmath}
\usepackage[colorlinks=true,urlcolor=blue]{hyperref}
\usepackage{titlesec,enumitem,flafter}
\usepackage{url}\Urlmuskip=0mu plus 1mu\relax
\pagestyle{plain}

% keep a heading with the text that follows it
\titleformat{\section}{\normalfont\large\bfseries}{\thesection.}{0.5em}{}
\titlespacing*{\section}{0pt}{1.4ex plus .3ex}{0.7ex}
\titlespacing*{\subsection}{0pt}{1.1ex plus .2ex}{0.5ex}
\widowpenalty=10000
\clubpenalty=10000
\brokenpenalty=10000
\setlength{\parskip}{0.35em}
% floats: prefer top of column, never mid-paragraph
\renewcommand{\topfraction}{0.95}
\renewcommand{\dbltopfraction}{0.95}
\renewcommand{\bottomfraction}{0.7}
\renewcommand{\textfraction}{0.08}
\renewcommand{\floatpagefraction}{0.6}
\renewcommand{\dblfloatpagefraction}{0.6}
\setcounter{topnumber}{2}
\setcounter{bottomnumber}{1}
\setcounter{totalnumber}{3}
\setlength{\textfloatsep}{10pt plus 2pt minus 2pt}
\setlength{\intextsep}{10pt plus 2pt minus 2pt}

\title{\vspace{-2.2em}\textbf{The 2026 Strait of Hormuz Closure:\\Crude Flows, Import Composition, and Retail\\Pass-Through in the United States and China}\vspace{-0.4em}}
\author{Henry Vianna \quad---\quad University of Toronto}
\date{\today}

\begin{document}
\maketitle
\vspace{-1.5em}

\section{RESEARCH QUESTION AND DESIGN}

How did the 2026 Strait of Hormuz closure change crude oil flows and
import composition for the US and China relative to a fixed pre-crisis
baseline, and how differently did the shock pass through to retail fuel
prices in each economy?

This comparison is motivated by an exposure asymmetry. The United States
imports around 0.43 mb/d of Hormuz-transiting crude, whereas China
imports around 5.4 mb/d through the same chokepoint (EIA, World Oil
Transit Chokepoints, 1H2025). This implies an exposure ratio between
China and the US of roughly 12:1. The two countries also differ
institutionally. US pump prices are determined by market fluctuations,
in contrast to China, where the NDRC administers prices under a
published formula. Therefore, crude pass-through to pump prices is
estimated econometrically for the US and measured as a policy wedge for
China. These are different quantities and are never placed in the same
regression.

\textbf{Design constraints.} Treatment is ongoing and multistage; there
is no clean post-period. There is no control group, since oil is a
single global market. The closure is confounded with the concurrent
Iran war, Iran's own export collapse, and the US naval blockade of
Iranian ports declared on 12 April. These cannot be separated using
observational data alone. Every residual reported below is therefore a
crisis-period effect rather than a closure-specific one.

Demand was deliberately excluded from the question, as it is not
observable at this horizon. The original interest in the effect on a
typical citizen was sharpened into retail price pass-through, which is
observable in both countries.

\textbf{Baseline.} The baseline is fixed at 28 Feb 2025 -- 27 Feb 2026
and is never rolling. Because most series are monthly, the window is
rendered two ways --- Mar-2025 to Feb-2026 and Feb-2025 to Jan-2026 ---
and every estimate is reported across both. The second window exists
because February 2026 shows evidence of pre-positioning, and including
an inflated month in the baseline would inflate every measured decline.

\section{DATA AND VALIDATION}

Seven processed datasets were built from EIA, JODI, GACC, FRED, NDRC and
WTO DataLab sources. The organising principle is one file per
observation frequency, and never a master table. Broadcasting monthly
data onto daily rows invents observations, so joins are performed at the
point of a specific question, which keeps the aggregation choice visible
in the analysis rather than buried in the data.

Several cleaning decisions materially affect the results. JODI reports
five units per observation, so failing to filter on \texttt{UNIT\_MEASURE}
quintuples every row, and \texttt{CONVBBL} is the tonne-to-barrel
conversion factor rather than a quantity. GACC publishes Table 14 in
three different column layouts across the thirteen files used, so rows
are located by commodity rather than by position. EIA's import series
places aggregate codes (\texttt{NUS-Z00}, \texttt{ME0}, \texttt{MN0},
\texttt{MP0}) alongside individual countries, so summing without
excluding them triple-counts.

\textbf{Validation.} Four independent checks were passed. US country
sums equal the \texttt{NUS-Z00} aggregate exactly in all 89 months.
Chinese volumes computed from GACC tonnages agree with JODI in 10 of 12
overlapping months. The SPR pipeline reproduces the announced 2022
strategic release at $-985.9$ kb/d against a stated 1 mb/d, using no
information about that policy. And the chain of announced NDRC price
changes reconciles with the published Beijing price levels across all 25
adjustment windows.

\section{DID THE DISRUPTION OCCUR?}

\begin{figure}[htbp]
\centering
\includegraphics[width=0.78\textwidth]{figures/hormuz-disruption-overview.png}
\caption{Strait of Hormuz transits, Mar 2025 -- Jul 2026. Shaded region
is the fixed baseline; dashed line marks the 28 Feb 2026 closure.}
\end{figure}

Tanker calls and cargo volumes both fell by a factor of 10 to 12 against
baseline. This was tested across three baseline definitions
(equal-length, pre-spike, and fixed-year) because a pre-closure run-up
in tanker calls --- above 80 per day in mid-February against 30 to 40 in
December and January --- suggested the baseline itself might be
contaminated by pre-positioning. The fixed-year baseline gives
$12.07\times$ and the other two give approximately $10\times$, so the
result does not depend on the choice.

Two further observations support the interpretation. Container traffic
collapsed alongside tanker traffic, which indicates the disruption was
not tanker-specific. And average cargo per transit rose from a stable
level near 47,000 tonnes per call to peaks above 120,000, meaning the
transits that continued were substantially fuller. This is consistent
with only escorted or high-value convoys crossing.

\section{FLOW DISRUPTION}

\subsection*{China}

Seasonal factors were estimated on 2019--2025 only, so that the
treatment period could not contaminate the profile. Each month was
normalised by its own year's mean, which removes level differences
between years, and the resulting ratios were collapsed by calendar
month. Chinese refinery intake proves close to aseasonal, with a range
of only 7 percentage points across the year and no spring turnaround
signature. Crude imports are more seasonal, with an 11-point range.
This matters because seasonality was the most obvious competing
explanation for the observed decline, and it can now be ruled out on
evidence rather than assertion.

\begin{table}[htbp]
\centering\small
\begin{tabular}{lrr}
\toprule
\textbf{2026} & \textbf{Intake} & \textbf{Imports} \\
\midrule
March & 87 (0.5\%) & 254 (2.1\%) \\
April & 2{,}412 (14.9\%) & 2{,}281 (19.2\%) \\
May & 4{,}138 (25.6\%) & 4{,}057 (34.2\%) \\
\bottomrule
\end{tabular}
\caption{Chinese seasonally adjusted shortfall below baseline, kb/d, with percentage of baseline in parentheses.}
\label{tab:china-flows}
\end{table}

March is
statistically indistinguishable from baseline. This is consistent with
the three to four week Gulf-to-China shipping lag, since March arrivals
were largely February loadings, and it means April is the first fully
treated month. The escalation profile across the three months matches
the event chronology.

The May import shortfall is approximately 75\% of China's 5.4 mb/d
Hormuz exposure. It is below the physical ceiling, as it must be, and
close enough to it to be consistent with the loss of most Hormuz supply
while non-Hormuz suppliers continued.

\textbf{No inventory buffering.} Crude entering Chinese refineries comes
from imports plus domestic production, net of stock movements, so the
gap between intake and imports isolates the latter two. That gap sat 24
and 12 kb/d above its baseline of 4,283 in April and May, which is under
half a percent. Refiners cut runs in proportion to the import shortfall
rather than drawing down inventories. Given that Chinese crude stocks
are commonly estimated near one billion barrels, this appears to have
been a choice rather than a constraint.

\subsection*{United States}

\begin{figure}[htbp]
\centering
\includegraphics[width=0.78\textwidth]{figures/us-quantities-bar.png}
\caption{US crude flow deviations from the Mar 2025 -- Feb 2026
baseline, kb/d. Not adjusted for seasonality.}
\end{figure}

Every monthly deviation in total US crude imports falls below $2\sigma$
against a month-over-month noise floor of 340 kb/d. This is not
evidence that no effect occurred. Maximum Hormuz exposure is 430 kb/d,
so a complete loss of every Hormuz barrel would register at only
$1.26\sigma$. In other words, the month-to-month
variation in total US imports is comparable in size to the entire
quantity at risk, so even the largest disruption physically possible
would be hidden inside ordinary fluctuation. The series cannot answer
the question, whatever the true effect turned out to be. This
calculation was performed before any estimation, so the standard was set
independently of the result.

The adjustment ran through other channels entirely. Measured against
series-specific noise floors:

\begin{table}[htbp]
\centering\small
\begin{tabular}{lrr}
\toprule
\textbf{May 2026} & \textbf{Deviation} & \textbf{$\sigma$} \\
\midrule
Crude exports & $+1{,}411$ & 4.2 \\
SPR draw & $-1{,}300$ & $>$10 \\
Commercial stock draw & $-855$ & --- \\
Crude imports & $-245$ & 0.7 \\
\bottomrule
\end{tabular}
\caption{US crude flow deviations for May 2026 against series-specific month-over-month noise floors.}
\label{tab:us-flows}
\end{table}

Exports rose by three to five times the amount by which
imports fell. If the United States had been scrambling to cover a
domestic shortfall, exports would have fallen. Instead, the strategic
reserve was drawn at a rate exceeding the 2022 release while more crude
was sent abroad. The US was not absorbing a supply shock so much as
supplying a tight global market, with the reserve financing it.

\section{IMPORT COMPOSITION}

\begin{figure}[htbp]
\centering
\includegraphics[width=0.78\textwidth]{figures/us-crude-import-composition.png}
\caption{US crude import shares by origin region. Dashed line marks the
closure.}
\end{figure}

Disaggregating by origin resolves what the aggregate could not.
Gulf-origin crude (Saudi Arabia, Iraq, Kuwait, UAE) fell 354 kb/d in May
2026, which is $2.3\sigma$ and a decline of roughly two thirds, and the
result holds across both baseline windows. The Gulf-origin series is
smaller and far less volatile than total imports, so the same disruption
registers as a much larger multiple of that series' own noise: $2.3\sigma$
here against $0.7\sigma$ in the aggregate. Measuring the channel that was
actually disrupted, rather than the total it sits inside, is what makes
the effect visible.

Latin America absorbed the entire shortfall and more, rising 420 kb/d at
$3.2\sigma$, while every other region fell. North America was flat at
$0.19\sigma$ despite supplying 4,248 kb/d at baseline. Canada is eight
times the Gulf's volume with obvious spare capacity, and it did not step
up.

This points to grade compatibility rather than volume availability as
the binding constraint. Gulf Coast refineries are configured for
medium-sour crude, and heavy-sour Latin American barrels are a closer
substitute than light-sweet Canadian synthetic. It should be stressed
that this is inferred from which regions moved, not from crude assay
data; contract structure or freight economics could produce the same
pattern.

February 2026 Gulf-origin imports ran 48\% above baseline before
collapsing. This is the same pre-positioning signature observed in
mid-February tanker calls and in the March commercial stock build.

\section{RETAIL PASS-THROUGH}

\begin{figure}[htbp]
\centering
\includegraphics[width=0.78\textwidth]{figures/passthrough-comparison.png}
\caption{Gasoline, actual against counterfactual. Left: Beijing maximum
retail price against the path implied by the NDRC formula. Right: US
retail against a projection from a pre-crisis pass-through model.}
\end{figure}

\subsection*{United States}

Augmented Dickey-Fuller tests fail to reject a unit root in levels for
all four series and reject decisively in first differences, so
estimation proceeds in differences. This is a necessary step rather
than a formality: a regression between two $I(1)$ series in levels is
spurious. Weekly changes in retail price are regressed on weekly
changes in crude at lags 0 through 8, with HAC standard errors
(maxlags = 8) to account for the serial correlation that overlapping lag
windows induce. The specification follows Borenstein, Cameron and
Gilbert (1997), simplified to a distributed lag. Daily crude is
aggregated to Friday-ending weeks and matched to the following Monday's
retail survey, so that no pump price is regressed on crude prices that
had not yet been observed.

Cumulative pass-through is the sum of the nine estimated coefficients
--- lag 0, the contemporaneous week, through lag 8 --- expressed in
dollars per gallon for each dollar per barrel of crude. Complete
mechanical pass-through is $1/42 = 0.0238$, since a barrel contains 42
gallons.

\begin{table}[htbp]
\centering\small
\begin{tabular}{llrr}
\toprule
\textbf{Product} & \textbf{Benchmark} & \textbf{Cumulative} & \textbf{s.e.} \\
\midrule
Gasoline & Brent & 0.0213 & 0.0029 \\
Gasoline & WTI & 0.0244 & 0.0038 \\
Diesel & Brent & 0.0253 & 0.0064 \\
Diesel & WTI & 0.0280 & 0.0068 \\
\bottomrule
\end{tabular}
\caption{Cumulative US pass-through, \$/gal per \$/bbl, estimated on pre-crisis data. Complete mechanical pass-through is $1/42 = 0.0238$.}
\label{tab:passthrough}
\end{table}

No $t$-statistic against the benchmark exceeds 1 in absolute
value. Pre-crisis US pass-through is therefore statistically
indistinguishable from complete, for both products and both benchmarks.
Brent and WTI also do not differ meaningfully from one another, which
means the near-tripling of the Brent-WTI spread during the crisis cannot
be read off these coefficients.

\textbf{Counterfactual.} Rather than re-estimating on the crisis window,
which would leave 22 observations against 10 parameters, the pre-crisis
model is applied to crisis-period crude with nothing fitted on crisis
data. Predicted weekly changes are cumulated from the last pre-crisis
retail level to give a counterfactual price path.

By late July, gasoline sat 77 cents per gallon and diesel 97 cents per
gallon above that path. However, the mean weekly residual is
insignificant, with $t$ between 0.94 and 1.42, and the three largest
weeks account for essentially the entire gap while the remaining
nineteen net to under half a cent. Those weeks are 9 March, 27 April
and 4 May, each immediately following an escalation event: the IRGC's
confirmation of the closure, the naval blockade, and the re-closure. In
the week of 4 May, retail rose 32.9 cents while crude was falling.

This is a discrete-shock pattern rather than a change in pass-through
regime. The natural mechanisms are refining margin expansion or a risk
premium at the pump, neither of which can be tested without crack spread
data.

\subsection*{China}

The NDRC publishes an implied price change only when it chooses not to
follow the formula. Substituting the actual change wherever no implied
figure exists therefore yields a complete full-pass-through series,
which is cumulated from the 3 February Beijing level to give a
counterfactual path.

Across the 24 consecutive adjustment windows from February 2025 to
January 2026, no narrowing notice was issued. Two were issued within
six weeks of the closure.

\begin{table}[htbp]
\centering\small
\begin{tabular}{lrrr}
\toprule
\textbf{Date} & \textbf{Implied} & \textbf{Actual} & \textbf{Absorbed} \\
\midrule
23 Mar 2026 & 2{,}205 & 1{,}160 & 1{,}045 \\
7 Apr 2026 & 800 & 420 & 380 \\
\bottomrule
\end{tabular}
\caption{NDRC interventions, yuan per tonne of gasoline. Implied figures are published by NDRC only where they deviate from the actual adjustment.}
\label{tab:ndrc}
\end{table}

In total 1,425 yuan per tonne of gasoline increase was
withheld, equivalent to \$24.42 per barrel, and 1,375 yuan per tonne
(\$26.89/bbl) for diesel. The conversion uses 8.5 bbl/tonne for gasoline
and 7.45 for diesel, reflecting their densities of roughly 0.74 and
0.845 t/m$^3$. The exchange rate is 6.865
CNY/USD, the DEXCHUS mean over March and April 2026, a period in which
the rate moved only 1.4\%. Beijing
gasoline peaked at 10,930 against a counterfactual 11.4\% higher, with
diesel 12.1\% higher.

The wedge steps twice and then holds perfectly flat across the remaining
eight adjustments, including a 950 yuan per tonne reduction on 3 July.
The constraint on prices was therefore active only while crude
was rising steeply, and was let go once it fell back. This says
something specific about \emph{when} an administered pricing regime
actually departs from its own formula, which a single average
pass-through figure would average away entirely.

It is worth being precise about the legal basis. Article 6 of the
Petroleum Price Management Measures already deducts the refiner
processing margin progressively above \$80 per barrel, and that
deduction is contained within the implied figure NDRC publishes. The
wedge is therefore Article 7 discretion in full, which requires State
Council approval. Article 7 is also cited routinely for sub-50 yuan
deferrals, so its invocation alone is not evidence of intervention; the
narrowing of a large adjustment is.

\section{SYNTHESIS}

The exposure asymmetry did not translate into a proportional consumer
outcome. China lost roughly a third of its crude imports and cut
refinery runs in step. The United States lost two thirds of a channel
worth 8\% of its imports and replaced it within a month. Yet US pump
prices rose 54\% between February and May, from \$2.91 to \$4.48 for
gasoline and \$3.72 to \$5.60 for diesel, while Beijing gasoline rose
26\%. The country with twelve times the physical exposure experienced
half the retail increase.

The explanation is institutional rather than physical. US pass-through
was complete before the crisis and remained so through it, so the crude
price reached consumers because nothing stood between them and it.
China's formula would have delivered a comparable increase, and the
state overrode it twice.

Products were hit harder than crude in both countries. Chinese refined
product imports fell 48.6\% year on year against 29.0\% for crude, and
in the US the diesel counterfactual gap of 97 cents exceeded gasoline's
77 cents, while in China diesel's absorbed wedge was larger per barrel
than gasoline's. Two countries, independent datasets, and the same
asymmetry. This is consistent with Hormuz crude being distillate-rich
and with the Gulf being a major refined product exporter in its own
right.

Each country used the buffer available to it. The US had exports, a
strategic reserve, commercial inventories and grade-flexible
substitutes, and drew on all four while allowing prices to clear. China
had administered prices and used them, while refinery throughput
absorbed the physical shortfall. Neither response is more effective in
general; they are different margins of adjustment available to
differently structured systems.

\section{CAVEATS AND LIMITATIONS}

\textbf{Identification.} There is no control group, and the closure
cannot be separated from the war, Iran's export collapse, or the US
blockade. Every residual reported is a crisis-period effect.

\textbf{Statistical power.} US crude import volume cannot detect this
shock at any plausible effect size. The counterfactual test rests on 22
observations and cannot rule out a modest persistent shift in
pass-through.

\textbf{Asymmetric design.} No free English-language source provides
Chinese crude imports by origin at monthly frequency, so composition is
analysed only for the US. Comtrade lacks Chinese monthly detail, and
GACC Table 16 is value-only at HS 2-digit, where Chapter 27 mixes crude
with coal, LNG and products, and omits Iraq, Kuwait and Angola.
Separately, China reports close to zero crude from Iran while Malaysia
shows \$1.83bn of Chapter 27 trade, so transshipment makes any
``imports from Iran'' series unusable.

\textbf{Origin is not route.} Saudi and UAE barrels can reach open water
via the Yanbu and Fujairah pipelines, so Gulf-origin US imports are an
upper bound on Hormuz-transiting volume, and part of the observed
decline may be rerouting rather than lost supply.

\textbf{Measurement error correlates with treatment.} AIS spoofing and
dark-vessel activity increased during the closure, biasing transit
counts in an unknown direction.

\textbf{Specification.} Working in first differences discards the
long-run equilibrium relationship an error-correction model would
capture, despite all four series being $I(1)$ at the same order.
Asymmetry between rising and falling crude prices is assumed away, even
though the window contains both. Kristoufek and Lunackova argue these
series are fractionally integrated, in which case the standard
$I(0)/I(1)$ dichotomy underlying the ADF test is misspecified;
fractional integration is not estimable on a sample of this length.

\textbf{Unverifiable inputs.} The NDRC crude basket weighting is not
published, so implied changes are taken as stated and cannot be
independently reproduced. Chinese inventory data is unpublished. The
Chinese pass-through fraction is endpoint-sensitive --- 0.62 measured at
the May peak and 0.47 at the end of July --- because a fixed absolute
wedge is divided by a moving denominator; the absolute figures are the
stable ones.

\textbf{Comparability.} NDRC standard products are 89-octane gasoline
and No. 0 diesel, which are not comparable in levels to US 87 AKI.
Changes are comparable; levels are not.

\textbf{Unresolved.} A JODI-GACC discrepancy of roughly 14,500 kbbl in
January and February 2026 offsets across the pair to 0.3\% but remains
undiagnosed, and February is the treatment-onset month. JODI reports
Chinese refining at 16.4 mb/d for January 2026 against EIA's 14.2 for
2024, which appears definitional but is unresolved. Kiel Policy Brief
206 models this exact closure and has not yet been reviewed against
these results.


\newpage
\section*{SOURCES AND REFERENCES}
\small
\setlength{\parskip}{0.3em}

\subsection*{Data}
\begin{description}[leftmargin=0.9em,style=nextline,itemsep=0.2em,parsep=0pt,topsep=0.2em]
\item[EIA API v2] US weekly petroleum supply (crude imports, exports, production,
refinery inputs, commercial and SPR stocks) and monthly crude imports by country of
origin. \url{https://www.eia.gov/opendata/}
\item[EIA, World Oil Transit Chokepoints] Country Analysis Brief.
\url{https://www.eia.gov/international/content/analysis/special_topics/World_Oil_Transit_Chokepoints/}
--- 1H2025 Hormuz flows of 20.9 mb/d total oil, of which approximately 15 mb/d crude
and condensate and 5.5 mb/d refined products, representing around 20\% of global
petroleum liquids consumption and one quarter of maritime traded oil. Exposure:
China 5.4 mb/d against US 0.43 mb/d. This paper uses the 20.9 mb/d total-liquids
figure for headline throughput.
\item[FRED, St.\ Louis Fed] \texttt{DCOILBRENTEU} and \texttt{DCOILWTICO}, daily
USD/bbl; \texttt{DEXCHUS}, daily CNY per USD. \url{https://fred.stlouisfed.org/}
\item[JODI Oil World Database] Chinese monthly refinery intake
(\texttt{REFINOBS}) and crude imports (\texttt{TOTIMPSB}) in KBD, with
\texttt{CONVBBL} = 7.32 bbl/tonne. \url{https://www.jodidata.org/}
\item[GACC] ``(14) Major Import Commodities in Quantity and Value'', monthly.
Quantity in 10,000 tonnes, value in US\$1,000.
\url{http://english.customs.gov.cn/statics/report/monthly.html} --- each month is
published at its own UUID URL with no constructable pattern; per-file URLs are
recorded in \texttt{docs/data\_inventory.md}.
\item[NDRC Price Department] Refined oil price adjustment notices, Jan 2025 --
Jul 2026. Standard products are 89-octane gasoline and No.~0 vehicle diesel, priced
in yuan per tonne as provincial and central-city maxima inclusive of consumption
tax, VAT and surcharges, effective 24:00 on the adjustment date.
\url{https://www.ndrc.gov.cn/xwdt/ztzl/gncpyjg/} --- per-notice URLs are recorded in
\texttt{china-retail-prices-events.csv}.
\item[WTO DataLab] Strait of Hormuz Trade Tracker; daily transits and cargo volumes.
\url{https://datalab.wto.org/Strait-of-Hormuz-Trade-Tracker}
\end{description}

\subsection*{Policy and Regulation}
\begin{description}[leftmargin=0.9em,style=nextline,itemsep=0.2em,parsep=0pt,topsep=0.2em]
\item[Petroleum Price Management Measures (\emph{Shiyou Jiage Guanli Banfa})]
NDRC, issued under Fagai Jiage [2016] No.~64, 13 January 2016.
\url{https://www.ndrc.gov.cn/xxgk/zcfb/tz/201601/t20160113_963566.html}
Article 6 sets the crude price bands: at or below \$40/bbl prices are set as if crude
were \$40; between \$40 and \$80 a normal processing margin applies; above \$80 the
refiner processing margin is progressively deducted toward zero; at or above \$130
prices are in principle not raised, or raised less, with accompanying fiscal
measures. Article 7 sets a ten-working-day adjustment cycle effective at 24:00 on
the publication date, provides that changes below 50 yuan/tonne are not made and
carry into the next cycle, and permits the NDRC, with State Council approval, to
suspend, delay or narrow an adjustment under special circumstances.
\item[US Treasury OFAC] General Licence X, 22 June 2026, authorising Iranian oil
sales through 21 August 2026.
\url{https://ofac.treasury.gov/recent-actions/20260622_33}
\end{description}

\subsection*{Literature}
\begin{description}[leftmargin=0.9em,style=nextline,itemsep=0.2em,parsep=0pt,topsep=0.2em]
\item[Bacon, R.\ W.\ (1991)] ``Rockets and Feathers: The Asymmetric Speed of
Adjustment of UK Retail Gasoline Prices to Cost Changes.''
\emph{Energy Economics} 13 (July).
\item[Borenstein, S., Cameron, A.\ C.\ \& Gilbert, R.\ (1997)] ``Do Gasoline Prices
Respond Asymmetrically to Crude Oil Price Changes?'' \emph{Quarterly Journal of
Economics} 112(1): 305--339. \texttt{doi:10.1162/003355397555118}.
\url{https://faculty.haas.berkeley.edu/borenste/download/QJE97GasAsym.pdf} ---
finds that nearly all of the response to a crude increase reaches the pump within
four weeks while decreases pass through gradually over eight. Methodological basis
for the US pass-through estimate here, simplified from an error-correction model to
a distributed lag in first differences.
\item[Kristoufek, L.\ \& Lunackova, P.] ``Rockets and Feathers Meet Joseph:
Reinvestigating the Oil-Gasoline Asymmetry on the International Markets.''
\texttt{arXiv:1407.5466} --- argues these series are fractionally integrated, so the
standard $I(0)/I(1)$ framework underlying error-correction specifications is
misspecified, and finds no asymmetry in six of seven national markets. Cited here
as a limitation; fractional integration is not estimable on this sample.
\item[FTC, Bureau of Economics] ``Gasoline Price Changes and the Petroleum
Industry.'' Survey of asymmetric pass-through evidence.
\item[Kiel Institute] Policy Brief 206, March 2026 --- models a Strait of Hormuz
closure scenario. Not yet reviewed against these results.
\end{description}

\subsection*{Chronology}
\begin{description}[leftmargin=0.9em,style=nextline,itemsep=0.2em,parsep=0pt,topsep=0.2em]
\item[UKOilWatch] Hormuz Timeline, 83 individually sourced events; primary index used
to construct \texttt{docs/chronology\_table.xlsx}.
\url{https://ukoilwatch.com/hormuz-timeline}
\item[Wikipedia] ``2026 Strait of Hormuz crisis''; ``2026 Strait of Hormuz campaign''.
\item[World Bank Data Blog] ``Strait of Hormuz disruption sends oil prices surging.''
\item[News reporting] Per-event URLs (CNBC, CBS News, CNN, Al Jazeera, Al-Monitor,
RFE/RL, gCaptain, Kpler via CNBC) are recorded in the Source column of
\texttt{chronology\_table.xlsx}. Most were transcribed from the UKOilWatch index
rather than retrieved individually; live-update pages in particular are overwritten
as stories develop and should be verified before quotation.
\end{description}

\subsection*{Conversion Factors}
7.32 bbl/tonne for Chinese crude (JODI's own \texttt{CONVBBL}); 8.5 bbl/tonne for
gasoline and 7.45 for diesel in the yuan-per-tonne to dollar-per-barrel conversion;
6.865 CNY/USD, the DEXCHUS mean over March--April 2026, over which the series ranged
1.4\%; 42 US gallons per barrel.

\end{document}